% BEGIN GENERATED ARTIFACT: def:elementary-conjunction-propositional-logic
\begin{definitionbox}{Definition (Elementary Conjunction)}
\begin{definition}[Elementary Conjunction]
\label{def:elementary-conjunction-propositional-logic}
An \emph{elementary conjunction} is a finite conjunction of
\hyperref[def:literal-propositional-logic]{literals}. That is, it is a
\hyperref[def:term-propositional-logic]{term} of the form
\[
L_1\land L_2\land\cdots\land L_k,
\]
where each \(L_i\) is a literal. The one-literal case is permitted.
\end{definition}
\end{definitionbox}

\begin{remark*}[Standard quantified statement]
\[
L_1\land L_2\land\cdots\land L_k,
\qquad
L_i\text{ is a literal for each }i.
\]
\end{remark*}

\begin{remark*}[Interpretation]
Elementary conjunctions are the conjunction terms used as the building blocks
of disjunctive normal form. They record simultaneous literal requirements.
\end{remark*}

\begin{example*}[Examples]
\[
P,\qquad P\land\neg Q,\qquad \neg P\land R\land S
\]
are elementary conjunctions.
\end{example*}

\begin{dependencies}
\begin{itemize}
  \item \hyperref[def:term-propositional-logic]{Term}
  \item \hyperref[def:literal-propositional-logic]{Literal}
\end{itemize}
\end{dependencies}
% END GENERATED ARTIFACT: def:elementary-conjunction-propositional-logic

% BEGIN GENERATED ARTIFACT: def:minterm-propositional-logic
\begin{definitionbox}{Definition (Minterm)}
\begin{definition}[Minterm]
\label{def:minterm-propositional-logic}
Let \(P_1,\ldots,P_n\) be a finite list of propositional variables. A
\emph{minterm} relative to \(P_1,\ldots,P_n\) is an
\hyperref[def:elementary-conjunction-propositional-logic]{elementary conjunction}
that contains exactly one of \(P_i\) or \(\neg P_i\) for each
\(i=1,\ldots,n\).
\end{definition}
\end{definitionbox}

\begin{remark*}[Standard quantified statement]
\[
M=L_1\land\cdots\land L_n,
\]
where for each \(i\), exactly one of \(P_i\) or \(\neg P_i\) appears among the
literals \(L_1,\ldots,L_n\).
\end{remark*}

\begin{remark*}[Interpretation]
A minterm relative to \(P_1,\ldots,P_n\) specifies one complete truth-value
assignment on those variables: choosing \(P_i\) requires \(P_i\) to be true,
and choosing \(\neg P_i\) requires \(P_i\) to be false.
\end{remark*}

\begin{example*}[Examples]
Relative to \(P,Q,R\), the formula
\[
P\land\neg Q\land R
\]
is a minterm. Relative to \(P,Q,R\), the formula \(P\land R\) is not a
minterm because it omits both \(Q\) and \(\neg Q\).
\end{example*}

\begin{dependencies}
\begin{itemize}
  \item \hyperref[def:elementary-conjunction-propositional-logic]{Elementary Conjunction}
  \item \hyperref[def:formula-variable-set-propositional-logic]{Variable Set of a Formula}
\end{itemize}
\end{dependencies}
% END GENERATED ARTIFACT: def:minterm-propositional-logic

% BEGIN GENERATED ARTIFACT: def:disjunctive-normal-form-propositional-logic
\begin{definitionbox}{Definition (Disjunctive Normal Form)}
\begin{definition}[Disjunctive Normal Form]
\label{def:disjunctive-normal-form-propositional-logic}
A formula is in \emph{disjunctive normal form}, abbreviated \emph{DNF}, when it
is a finite disjunction of
\hyperref[def:elementary-conjunction-propositional-logic]{elementary conjunctions}.
Equivalently, it has the shape
\[
C_1\lor C_2\lor\cdots\lor C_m,
\]
where each \(C_j\) is an elementary conjunction.
\end{definition}
\end{definitionbox}

\begin{remark*}[Standard quantified statement]
\[
\varphi\text{ is in DNF}
\Longleftrightarrow
\varphi=C_1\lor C_2\lor\cdots\lor C_m
\]
for finitely many elementary conjunctions \(C_1,\ldots,C_m\).
\end{remark*}

\begin{remark*}[Interpretation]
DNF is a syntactic form, but it is used semantically: once a formula is replaced
by a \hyperref[def:logical-equivalence-propositional-logic]{logically equivalent}
DNF formula, questions about
\hyperref[def:satisfiable-formula-propositional-logic]{satisfiability} reduce
to questions about individual conjunction terms.
\end{remark*}

\begin{example*}[Examples]
\[
(P\land Q)\lor(\neg P\land R)\lor S
\]
is in DNF because each disjunct is an elementary conjunction.
\end{example*}

\begin{dependencies}
\begin{itemize}
  \item \hyperref[def:elementary-conjunction-propositional-logic]{Elementary Conjunction}
  \item \hyperref[def:logical-equivalence-propositional-logic]{Logical Equivalence}
  \item \hyperref[def:satisfiable-formula-propositional-logic]{Satisfiable Formula}
\end{itemize}
\end{dependencies}
% END GENERATED ARTIFACT: def:disjunctive-normal-form-propositional-logic

% BEGIN GENERATED ARTIFACT: thm:existence-disjunctive-normal-form-propositional-logic
\begin{theorembox}{Theorem (Existence of Disjunctive Normal Form)}
\begin{theorem}[Existence of Disjunctive Normal Form]
\label{thm:existence-disjunctive-normal-form-propositional-logic}
For every formula \(\varphi\in\WFF\), there exists a formula
\(\psi\in\WFF\) such that \(\psi\) is in
\hyperref[def:disjunctive-normal-form-propositional-logic]{disjunctive normal form}
and
\[
\varphi\equiv\psi.
\]
\noindent\hyperref[prf:existence-disjunctive-normal-form-propositional-logic]{Go to proof.}
\end{theorem}
\end{theorembox}

\begin{remark*}[Standard quantified statement]
\[
(\forall\varphi\in\WFF)(\exists\psi\in\WFF)
\bigl(\psi\text{ is in DNF}\land \varphi\equiv\psi\bigr).
\]
\end{remark*}

\begin{remark*}[Interpretation]
One first converts \(\varphi\) to negation normal form, then repeatedly applies
distributive laws to move conjunctions below disjunctions. Replacement of
logical equivalents preserves equivalence at each step.
\end{remark*}

\begin{dependencies}
\begin{itemize}
  \item \hyperref[def:disjunctive-normal-form-propositional-logic]{Disjunctive Normal Form}
  \item \hyperref[def:negation-normal-form-propositional-logic]{Negation Normal Form}
  \item \hyperref[thm:distributive-laws-propositional-logic]{Distributive Laws for Propositional Logic}
  \item \hyperref[thm:replacement-of-logical-equivalents-propositional-logic]{Replacement of Logical Equivalents}
  \item \hyperref[thm:existence-negation-normal-form-propositional-logic]{Existence of Negation Normal Form}
\end{itemize}
\end{dependencies}
% END GENERATED ARTIFACT: thm:existence-disjunctive-normal-form-propositional-logic

% BEGIN GENERATED ARTIFACT: thm:dnf-satisfiability-criterion-propositional-logic
\begin{theorembox}{Theorem (DNF Satisfiability Criterion)}
\begin{theorem}[DNF Satisfiability Criterion]
\label{thm:dnf-satisfiability-criterion-propositional-logic}
Let \(\psi\) be a formula in
\hyperref[def:disjunctive-normal-form-propositional-logic]{disjunctive normal form}.
Then \(\psi\) is
\hyperref[def:satisfiable-formula-propositional-logic]{satisfiable} if and only
if at least one conjunction term in \(\psi\) contains no
\hyperref[def:complementary-literals-propositional-logic]{complementary pair of literals}.
\noindent\hyperref[prf:dnf-satisfiability-criterion-propositional-logic]{Go to proof.}
\end{theorem}
\end{theorembox}

\begin{remark*}[Standard quantified statement]
\[
\psi\text{ is a satisfiable DNF formula}
\Longleftrightarrow
\text{some conjunction term of }\psi\text{ has no complementary literals.}
\]
\end{remark*}

\begin{remark*}[Interpretation]
A DNF formula is true exactly when at least one of its conjunction terms is
true. A conjunction term can be made true exactly when it does not require both
a variable and its negation.
\end{remark*}

\begin{dependencies}
\begin{itemize}
  \item \hyperref[def:disjunctive-normal-form-propositional-logic]{Disjunctive Normal Form}
  \item \hyperref[def:complementary-literals-propositional-logic]{Complementary Literals}
  \item \hyperref[def:satisfiable-formula-propositional-logic]{Satisfiable Formula}
\end{itemize}
\end{dependencies}
% END GENERATED ARTIFACT: thm:dnf-satisfiability-criterion-propositional-logic
